\documentclass[11pt,a4paper]{article}
\usepackage[spanish,es-nodecimaldot,es-noquoting]{babel}
\usepackage{fontspec}
\setmainfont{texgyrepagella-regular.otf}[
 BoldFont=texgyrepagella-bold.otf,ItalicFont=texgyrepagella-italic.otf,
 BoldItalicFont=texgyrepagella-bolditalic.otf]
\setsansfont{texgyreheros-regular.otf}[BoldFont=texgyreheros-bold.otf,
 ItalicFont=texgyreheros-italic.otf]
\usepackage{amsmath,amssymb,amsthm,mathtools}
\usepackage{unicode-math}
\setmathfont{texgyrepagella-math.otf}
\usepackage[a4paper,top=25mm,bottom=25mm,left=28mm,right=28mm]{geometry}
\usepackage{microtype,booktabs,tabularx,array,enumitem,needspace,etoolbox,placeins}
\usepackage{tikz}
\usetikzlibrary{arrows.meta,positioning,calc}
\usepackage{xcolor}
\definecolor{ink}{HTML}{183D4A}
\definecolor{accent}{HTML}{8C3D2E}
\usepackage[colorlinks=true,linkcolor=ink,citecolor=ink,urlcolor=ink]{hyperref}
\usepackage{fancyhdr}
\pagestyle{fancy}
\fancyhf{}
\fancyhead[L]{\small\itshape Holografía modular triádica}
\fancyhead[R]{\small Acción, geometría e información}
\fancyfoot[C]{\thepage}
\setlength{\headheight}{14pt}
\setlength{\footskip}{13mm}
\pretocmd{\section}{\par\Needspace{12\baselineskip}}{}{}
\pretocmd{\subsection}{\par\Needspace{9\baselineskip}}{}{}
\raggedbottom
\setlength{\parindent}{1.2em}
\setlength{\parskip}{3pt}
\setlength{\emergencystretch}{2em}
\clubpenalty=10000
\widowpenalty=10000
\displaywidowpenalty=10000
\setlist{nosep,leftmargin=1.5em}
\numberwithin{equation}{section}
\newtheorem{theorem}{Teorema}[section]
\newtheorem{proposition}[theorem]{Proposición}
\newtheorem{lemma}[theorem]{Lema}
\theoremstyle{definition}
\newtheorem{definition}[theorem]{Definición}
\AtBeginEnvironment{definition}{\Needspace{8\baselineskip}}
\AtBeginEnvironment{proposition}{\Needspace{8\baselineskip}}
\AtBeginEnvironment{theorem}{\Needspace{8\baselineskip}}
\newcommand{\HMT}{\mathrm{HMT}}
\newcommand{\Tr}{\operatorname{Tr}}
\newcommand{\Spec}{\operatorname{Spec}}
\newcommand{\Cstar}{C^{*}}
\newcommand{\dd}{\mathrm d}
\newcommand{\R}{\mathbb R}
\newcommand{\Z}{\mathbb Z}
\newcommand{\I}{\mathrm I}
\newenvironment{apunte}[1]{\par\Needspace{6\baselineskip}\smallskip
\begin{quote}\small\raggedright\color{ink}\noindent\textbf{Apunte de borrador: #1.}\ }
{\end{quote}\smallskip}
\hypersetup{pdftitle={Acción, geometría e información en Holografía Modular Triádica},pdfauthor={Material de trabajo dirigido por Rubén Ramos Balsa},pdfsubject={Prototipo expositivo; fuentes HMT–MD y desarrollos de septiembre de 2026}}
\title{\color{ink}Acción, geometría e información en Holografía Modular Triádica\\[5pt]\large El funcional de Barbero--Immirzi y la estructura de área}
\author{Rubén Ramos Balsa}
\date{\small 6 de septiembre de 2026; revisión de 7 de septiembre\\[3pt]
\textcolor{accent}{Borrador 0.2 para discusión con el autor}}
\begin{document}
\maketitle
\begin{abstract}
Se desarrolla la relación entre la estructura generadora de Holografía
Modular Triádica, sus realizaciones cuadráticas y la construcción de la
acción. La exposición parte de las posiciones y operaciones de APP,
introduce los regímenes orientados del TRIT y mantiene la memoria del
transporte TPK. El sello dodecafásico se caracteriza mediante su reconstrucción
integral y su lectura periódica reversible. A partir de las publicaciones
numéricas y angulares se estudia el funcional de Barbero--Immirzi, cuyo
contenido de incidencia determina la normalización del área.
La lectura conjunta del área y del Hamiltoniano modular distingue las
ocupaciones de borde que una lectura total identifica. La reciprocidad
entre las secciones de acción y gravedad conserva, en la realización
circular declarada, el área de Planck.
La métrica, la conexión con torsión y el Hamiltoniano modular se presentan
con sus funciones respectivas. Este borrador reúne las demostraciones
locales y desarrolla su interpretación estructural; las notas editoriales
identifican las integraciones necesarias para la autonomía del artículo.
\end{abstract}\noindent\textbf{Palabras clave:} Holografía Modular Triádica; acción; Barbero--Immirzi;
memoria de transporte; geometría areal; holonomía.
\medskip

\noindent\small\textit{Nota editorial.} Este documento es un material preparatorio
para discusión del argumento, de la notación y de la selección de pruebas.
No reemplaza el tratado ni constituye una entrega científica final.
Las afirmaciones de cada resultado se circunscriben a sus dominios explícitos.
La asistencia de inteligencia artificial se ha utilizado en la organización,
redacción y comprobación local; la autoría de las fuentes se conserva en el
expediente de procedencia.\normalsize
\clearpage

\section{Fundamentos discretos y organización geométrica}
\label{sec:inicio}
\subsection{Marcas interiores y disposición transversal}
Se parte del camino marcado
\[
\mathsf P_{10}:0\longrightarrow1\longrightarrow\cdots
\longrightarrow9\longrightarrow10.
\]
Los extremos delimitan la carta y las nueve posiciones interiores forman
$\mathcal D_9=\{1,\ldots,9\}$. La representación mediante un segmento
sirve para ordenar marcas: en este punto no se presupone una recta real
completa. La fuente distingue expresamente el soporte finito y su
publicación límite posterior \cite{segmento}.

\begin{definition}[Carta positiva y continuación con acarreo]
La carta $\lambda_9:\mathcal D_9\to\Z/9\Z$ envía $a$ a su clase
residual, con $\lambda_9(9)=[0]_9$. El sucesor visible $u_+$ y su
incremento de memoria $\kappa_+$ se definen por
\[
u_+(a)=\begin{cases}a+1,&a<9,\\1,&a=9,\end{cases}
\qquad
\kappa_+(a)=\begin{cases}0,&a<9,\\1,&a=9.\end{cases}
\]
Se cumple exactamente $a+1=u_+(a)+9\kappa_+(a)$.
\end{definition}
El retorno $9\to1$ afecta al representante; el acarreo distingue
la prolongación del estado. Para cualquier entero positivo $n$,
\[
n=9q_9(n)+\operatorname{dr}_9(n),\quad
q_9(n)=\left\lfloor\frac{n-1}{9}\right\rfloor,\quad
\operatorname{dr}_9(n)=1+((n-1)\bmod9).
\]
La segunda disposición, transversal a la primera, construye
$\mathcal D_9\times\mathcal D_9$. Sus coordenadas se representan
horizontal y verticalmente, antes de identificar los lados mediante
la continuación modular. El resultado es la carta de 81 celdas
del toro aritmético $X_9$.

\begin{figure}[htbp]
\centering
\begin{tikzpicture}[x=6.8mm,y=6.8mm,>=Latex,font=\small]
\fill[ink!5] (.5,.5) rectangle (9.5,9.5);
\draw[step=1,ink!15] (.5,.5) grid (9.5,9.5);
\draw[->,ink,thick] (0,0)--(10.6,0) node[right] {$j$};
\draw[->,ink,thick] (0,0)--(0,10.6) node[above] {$i$};
\foreach \k in {0,...,10}{
 \draw (\k,.09)--(\k,-.09) node[below] {\k};
 \draw (.09,\k)--(-.09,\k);
}
\foreach \k in {1,...,10}{\node[left] at (-.12,\k) {\k};}
\foreach \i in {1,...,9}{\foreach \j in {1,...,9}{
 \pgfmathtruncatemacro{\v}{1+mod(\i*\j-1,9)}
 \node[font=\footnotesize] at (\j,\i) {\v};
}}
\draw[accent,thick] (.5,8.5) rectangle (9.5,9.5);
\draw[accent,thick] (8.5,.5) rectangle (9.5,9.5);
\node[align=left,font=\small,text width=43mm,anchor=west] at (11,6.5)
 {Una posición\\dos evaluaciones\\$\Sigma(i,j)$ y $\Pi(i,j)$};
\node[align=left,font=\small,text width=43mm,anchor=west] at (11,3.3)
 {Tabla mostrada:\\$\Pi(i,j)=\operatorname{dr}_9(ij)$\\Borde marcado: clase nula};
\end{tikzpicture}
\caption{Construcción horizontal y vertical de la carta APP. Se dibuja el
representante multiplicativo de cada celda. El orden ascendente de la
coordenada vertical es una convención gráfica; las rutas cardinales se
declaran después. Las posiciones y sus valores aritméticos son objetos distintos.}
\label{fig:app-construction}
\end{figure}

Cada fila multiplicativa es una acumulación del mismo incremento:
$ij=i+\cdots+i$ con $j$ sumandos, y cada columna expresa la
acumulación transversal. Las evaluaciones suma y producto actúan en el
mismo soporte; no constituyen dos APP independientes.
El gnomon de la última fila y columna tiene $9+8=17$ celdas,
de modo que $81=64+17$. La extensión volumétrica distingue
$1000=729+243+27+1$, preservando interior, caras, aristas y cierre.
No se identifican esos estratos únicamente porque alguno comparta cardinal
con un espacio de palabras.

\begin{apunte}{orden compartido de los dos artículos}
La apertura coordinada con el artículo de $\alpha$ y del electrón será:
segmento marcado, carta vertical, operaciones y caminos, semillas,
orientación TRIT y composición TPK. Este dibujo y sus definiciones se
proponen para ambos borradores; la extensión expositiva podrá diferir.
Los nueve lugares interiores no se confunden con los diez intervalos
consecutivos que separan las once marcas de $\mathsf P_{10}$.
\end{apunte}

\subsection{La acción como estructura y como lectura}
La pregunta de este artículo no es únicamente qué valor publica una constante,
sino qué operaciones producen ese valor y qué información permite recuperarlas.
En Holografía Modular Triádica, abreviada HMT, un registro conserva la hoja
aritmética, la orientación, el cociente, el acarreo y la historia de transporte.
La evaluación escalar es una aplicación de ese registro; no lo sustituye.

El núcleo de partida está constituido por APP, TRIT y el
\emph{Topological Phase Kernel} (TPK). APP proporciona el soporte y sus
operaciones; TRIT determina la orientación local; TPK compone los transportes,
la memoria y la prolongación compatible. Las publicaciones
$\pi_{\HMT},\varphi_{\HMT},e_{\HMT},\alpha_{\HMT}$ pertenecen a una misma
genealogía, con un orden constructivo interno. Las primeras lecturas pueden
preceder a la clausura dodecafásica de $\alpha_{\HMT}$ sin convertirse
en valores convencionales introducidos desde fuera.

El paso hacia Mecánica Dimensional se efectúa mediante los lectores de
acción, duración y magnitud. Este artículo se concentra en una composición
particular: el registro produce una escala y dos secciones orientadas de acción;
la rama angular publica los canales conjugados; la incidencia de borde
determina un funcional cuya realización se identifica con el parámetro de
Barbero--Immirzi. La geometría excepcional procede de la proyección de
incidencia del mismo estado. La coincidencia de origen no identifica sus
codominios: se conservan los mapas entre ellos.

\begin{figure}[htbp]
\centering
\begin{tikzpicture}[>=Latex,node distance=10mm,
box/.style={draw=ink,rounded corners=2pt,align=center,font=\small,
text width=40mm,minimum height=12mm,inner sep=4pt},
every path/.style={draw=ink}]
\node[box,text width=134mm] (base) at (0,0) {APP--TRIT--TPK\\Estado, orientación, acarreo, memoria y prolongación compatible};
\node[box] (num) at (-3.6,-1.8) {Lecturas numéricas\\$\pi,\varphi,e,\alpha$};
\node[box] (exc) at (3.6,-1.8) {Incidencia excepcional\\Witt, Golay, Leech};
\node[box] (act) at (-3.6,-3.6) {Recta de acción\\$H_5,\eta_{\rm ret},\Sigma_H$};
\node[box] (inc) at (3.6,-3.6) {Borde marcado\\Incidencias $90/120$};
\node[box,text width=52mm] (gam) at (0,-5.4) {Funcional angular y areal\\$\gamma_{\HMT}$};
\draw[->] (base.south)-|(num.north);
\draw[->] (base.south)-|(exc.north);
\draw[->] (num)--(act);
\draw[->] (exc)--(inc);
\draw[->] (act.south)|-(gam.west);
\draw[->] (inc.south)|-(gam.east);
\end{tikzpicture}
\caption{Organización de las dependencias del artículo. Las flechas distinguen
publicaciones y transportes; no representan una identificación entre un valor
numérico y un grupo excepcional. La construcción angular se explicita en la sección~\ref{sec:action}.}
\end{figure}
\FloatBarrier

\subsection{Soporte y memoria aritmética}
En la fuente APP, el soporte es $X_9=(\Z/9\Z)^2$, con caminos cardinales.
La raíz digital positiva se escribe
\[
\operatorname{dr}(n)=1+((n-1)\bmod9).
\]
Las dos evaluaciones de un vértice $(i,j)$ son
$\Sigma(i,j)=\operatorname{dr}(i+j)$ y
$\Pi(i,j)=\operatorname{dr}(ij)$. Sus elevaciones satisfacen
\[
i+j=R^+_{ij}+9q^+_{ij},
\qquad
ij=R^\times_{ij}+9q^\times_{ij}.
\]
Los cocientes permiten recuperar lo que la reducción oculta. En una
transición ternaria balanceada, $a+b=r+3c$ distingue igualmente el residuo
$r$ y el acarreo $c$. Los estados $(0,1),(0,0),(0,-1)$ tienen una misma
lectura visible y distinta información de transporte.

\begin{definition}[Semilla orientada y trayectoria]
Sea $\mathscr D=\{N,E,S,O\}$ el alfabeto cardinal. Una condición inicial
de doble cursor pertenece a
$(X_9\times\mathscr D)^2$; su cardinal es
$(81\cdot4)^2=104\,976$. Una semilla no es la palabra que publica:
el censo de la fuente distingue 468 emisiones decimales $U_6$,
243 palabras ternarias visibles $w_6$ y el ambiente de 729 palabras
de longitud seis. Las fibras conservan la multiplicidad de procedencia.
\end{definition}
Para $x_0\in X_9$ y $w=d_1\cdots d_r$, la ruta es
$x_k=x_0+\sum_{j=1}^k\delta(d_j)\pmod9$.
El transporte cardinal conserva el punto inicial y la palabra.
El censo no sustituye la prolongación: ésta compone los niveles
$w_6,w_{12},w_{18},w_{24},w_{30},R_{36},G_9$,
con retorno de fase y avance de memoria.

\begin{apunte}{autonomía del artículo}
La versión individual debe explicar el censo y la prolongación conjunta,
no pedir al lector que acepte 468 o 243 como datos sin antecedente.
La recopilación conserva la fuente completa; aquí se fijan los tipos
para desarrollar esa prueba en la siguiente revisión sin alterar las
construcciones posteriores de acción.
\end{apunte}

El número genealógico se organiza mediante secciones compatibles de los
cilindros supervivientes. La contracción de esos cilindros publica la
realización arquimediana; la proyección conjugada conserva incidencia.
Las cinco construcciones del continuo permanecen como componentes del mismo
objeto, no como antecedentes escogidos por el resultado físico buscado.
La exposición completa de este núcleo y sus figuras se comparte
documentalmente con el artículo dedicado a $\alpha$ y al electrón
\cite{app,trit,tpk,witt}. El presente prototipo desarrolla a continuación
las consecuencias de registro y acción que le son propias.

\input{sections/geometria.tex}

\section{Actualización dodecafásica y memoria publicada}
\label{sec:eta}
Sea $d_t$ la dirección producida por una traza TPK. Las doce ventanas
$E_m=\{9m+1,\ldots,9m+9\}$, $0\le m<12$, dividen los 108 pasos del registro.
La orientación es
\[
\Sigma_{12}=(+,+,+,-,-,-,+,+,+,-,-,-).
\]
El evento de una ventana es la salida
\begin{equation}
a_m=\Sigma_{12}(m)
\sum_{t\in E_m}\mathbf1_{\{2,4,6,8\}}(d_t).
\label{eq:events}
\end{equation}
La notación $a_m$ evita confundir este evento con el número de Euler.
El registro fuente conserva además la orientación TRIT, la firma,
el acarreo y el segmento de historia correspondiente \cite{registro}.

\begin{proposition}[Incremento producido por el transporte]
Sea $\mathbf e_m$ la base de $\R^{12}$ y
$S\mathbf e_m=\mathbf e_{m+1\pmod{12}}$. Si
$z_{m+1}=z_m+a_m\mathbf e_m$ y $y_m=S^{-m}z_m$, entonces
\begin{equation}
y_{m+1}=S^{-1}y_m+\eta_m,
\qquad
\eta_m=a_mS^{-1}\mathbf e_0.
\label{eq:eta}
\end{equation}
Después de una vuelta completa,
\begin{equation}
y_{12}=y_0+\sum_{m=0}^{11}a_m\mathbf e_m.
\label{eq:cycle}
\end{equation}
\end{proposition}
\begin{proof}
Multiplicar la actualización de $z_m$ por $S^{-(m+1)}$ da
\[
y_{m+1}=S^{-1}y_m+a_mS^{-(m+1)}\mathbf e_m
=S^{-1}y_m+a_mS^{-1}\mathbf e_0.
\]
Al iterar, el coeficiente del evento $a_m$ es
$S^{-(12-m)}\mathbf e_0=\mathbf e_m$. Como $S^{12}=I$,
se obtiene \eqref{eq:cycle}.
\end{proof}

Por tanto, $\eta_m$ no es una fuerza libre elegida para imponer una
conservación: es una publicación del evento TPK. La igualdad de fase
después de una vuelta no cancela los eventos acumulados. Este lector
conserva los recuentos, mientras el registro completo conserva además su
orden y procedencia. Se distingue $\eta_m$, incremento vectorial, de
$\eta_{\rm ret}$, mantisa escalar de acción, y de $\eta_{\rm el}$,
parámetro de forma que se introducirá en la sección~\ref{sec:ellipse}.

\subsection{Balance cuadrático}
Definimos
\[
D_3=I-S^3,\quad D_4=I-S^4,\quad
B=D_3^{\mathsf T}D_3+D_4^{\mathsf T}D_4,\qquad
\mathcal M(y)=\tfrac12 y^{\mathsf T}By.
\]
Se cumplen $S^{\mathsf T}BS=B$ y $B_{00}=4$. Sustituyendo
\eqref{eq:eta} se deduce exactamente
\begin{equation}
\mathcal M(y_{m+1})-\mathcal M(y_m)
=a_m(By_m)_0+2a_m^2.
\end{equation}
Asimismo, la carga $q(y)=\sum_jy_j$ satisface
$q(y_{m+1})-q(y_m)=a_m$.
La prueba es la expansión del polinomio cuadrático tras retirar la
permutación $S^{-1}$. Ningún valor metrológico entra en ese cálculo.
La forma $\mathcal M$ mide este registro; su interpretación como una
energía física determinada requeriría el lector dimensional correspondiente.

\input{sections/sello_y_lectura.tex}

\section{Descomposición cíclica y representaciones determinantales}
Los proyectores cíclicos
\[
P_3=\tfrac14(I+S^3+S^6+S^9),\qquad
P_4=\tfrac13(I+S^4+S^8)
\]
tienen imágenes de dimensiones tres y cuatro. Si $E=P_3-P_4$, entonces
\begin{equation}
\frac{\Tr E}{12}=-\frac1{12}.
\label{eq:minus12}
\end{equation}
El escalar procede de una diferencia de rangos normalizada, no de una
suma ordinaria de enteros positivos.

\begin{theorem}[Representación determinantal del lector de canal]
Sean $\mathcal H_3=\operatorname{Fix}(S^3)$ y
$\mathcal H_4=\operatorname{Fix}(S^4)$. Para $0<s<1$,
\begin{equation}
F(s)=\frac{\det_{\mathcal H_3}(I-sS)}{\det_{\mathcal H_4}(I-sS)}
=\frac{1-s^3}{1-s^4}
=\frac{1+s+s^2}{1+s+s^2+s^3}.
\label{eq:det}
\end{equation}
Su prolongación continua en $s=1$ vale $3/4$.
\end{theorem}
\begin{proof}
La restricción de $S$ a cada imagen es un ciclo de tres o cuatro
posiciones, respectivamente. Sus polinomios característicos dan
$\det(I-sS)=1-s^r$ para $r=3,4$. Al cancelar el factor común $1-s$,
queda la última expresión, regular en $s=1$.
\end{proof}

La sustitución $s=q^{30}$ publica
$F(q^{30})=(1-q^{90})/(1-q^{120})$.
El mismo par de espacios que produce \eqref{eq:minus12}
produce el cociente de canal. No se identifica por esta igualdad
una matriz de doce coordenadas con todo el estado TPK.

\subsection{Información orientada y sector complejo}
En el sector antipodal $V_-=\operatorname{Ran}((I-S^6)/2)$,
el operador $J=S^3$ satisface $J^2=-I$ y $J^{\mathsf T}=-J$.
Con $Q=P_4(I-S^6)/2$ se tiene
\[
\operatorname{rank}Q=2,\qquad E|_{V_-}=-Q,\qquad B|_{V_-}=5I-3Q.
\]
En $\operatorname{Ran}Q$, $S^2=-I$ y el determinante es $1+s^2$.
Así,
\[
F(s)=\frac{1+s+s^2}{1+s}\,\frac1{1+s^2}.
\]
La inversión de orientación cambia el signo de la forma
$\omega(u,v)=\langle Ju,v\rangle$, pero conserva $B$ y $F$.
La evaluación escalar no distingue todos los transportes orientados.

La traza transportada produce además
\[
a_n^{\rm tr}=\Tr(S^nE)=3\mathbf1_{3\mid n}-4\mathbf1_{4\mid n},
\]
y, para $|s|<1$,
\[
-\log F(s)=\sum_{n\ge1}\frac{a_n^{\rm tr}}n\,s^n.
\]
Esta familia de momentos conecta el presente bloque con los lectores
espectrales recopilados. Su desarrollo aritmético completo, incluido
Catalán, tendrá una residencia propia; no se presupone aquí que una
identidad de funciones decida por sí sola su naturaleza aritmética.

\section{Construcción de la escala y de las secciones de acción}
\label{sec:action}
Esta rama utiliza $\alpha_{\HMT}$ y $\varphi_{\HMT}$ ya publicadas.
En el cierre local de cuatro hojas aparece
\[
H_4=
\begin{pmatrix}
1&1&1&1\\1&1&-1&-1\\1&-1&1&-1\\1&-1&-1&1
\end{pmatrix}.
\]
Los divisores determinantes son $1,2,4,16$; por tanto,
\[
\operatorname{SNF}(H_4)=\operatorname{diag}(1,2,2,4),
\qquad
G_H=\operatorname{coker}H_4\simeq
\Z/2\Z\oplus\Z/2\Z\oplus\Z/4\Z.
\]
El lector de norma de la fuente aplica la sección constante
$s_\alpha(g)=\alpha_{\HMT}$ a los dieciséis elementos:
\begin{equation}
\Sigma_H=\varphi_{\HMT}
\prod_{g\in G_H}s_\alpha(g)
=\varphi_{\HMT}\alpha_{\HMT}^{16}.
\end{equation}
La fórmula distingue el índice del conúcleo, la norma sobre él y la
actuación única de la autoescala. El texto propietario desarrolla el
lector y sus cotas \cite{accion}. La valoración decimal publicada satisface
\[
\alpha_{\HMT}^{16}<10^{-34}
<\Sigma_H<10^{-33},
\qquad \lfloor\log_{10}\Sigma_H\rfloor=-34.
\]
El exponente así obtenido pertenece a la lectura interna. La carta
de unidades físicas es una operación posterior.

\subsection{Carácter regional, retorno orientado y coordenadas angulares}
La rama regional produce \cite{retorno}
\begin{align}
H_5={}&\frac12\sqrt{\frac{1000\alpha}{\varphi}}-\alpha
+\frac9{16}\alpha^2-\frac59\alpha^3
+\frac7{48}\alpha^4-\frac1{54}\alpha^5,\\
x_A={}&\frac{50\pi}{9}\alpha,\qquad D_A=e^{-2x_A},\\
\mathscr C_\pi(\alpha)={}&169\alpha^6+
\frac{D_A\alpha^7}{1-D_A\alpha},\\
\eta_{\rm ret}={}&H_5-\frac{90}{\pi}\mathscr C_\pi(\alpha).
\end{align}
En estas ecuaciones $\pi,\varphi,e,\alpha$ designan las publicaciones
HMT. Los coeficientes y el selector $169$ conservan sus propietarios
regionales en la recopilación; este prototipo no sustituye su construcción
por la mera exhibición de las fórmulas.

Sea $\mathcal L_S$ la recta de acción con base positiva $\mathcal U_S$.
Las secciones son
\begin{equation}
\hbar_{\rm pre}=H_5\,10^{-34}\mathcal U_S,\qquad
\hbar_{\rm ret}=\eta_{\rm ret}\,10^{-34}\mathcal U_S.
\end{equation}
Sus lecturas aditiva y multiplicativa,
\[
\delta_{2w}=H_5-\eta_{\rm ret},\qquad
R_{\rm act}=\frac{H_5}{\eta_{\rm ret}},
\]
conservan aspectos distintos del retorno. Para $h_a=2\pi\hbar_a$,
la acción de una fase obedece
\[
\mathcal S(\theta)=\hbar_a\theta_{\rm rad}
=h_a\frac{\theta_{\deg}}{360}.
\]
La rama angular se publica después en sus propias coordenadas:
\[
A_{\deg}=1000\alpha,\qquad
C^*_{\deg}=2(\eta_{\rm ret}+\alpha),\qquad
q_\pm=\exp\left[-\frac{\pi}{180}(A\pm C^*)\right].
\]
La distinción entre el incremento $\eta_m$ de \eqref{eq:eta}
y la mantisa $\eta_{\rm ret}$ es esencial: no se identifican
por compartir una letra ni mediante una analogía de retorno.

\section{Incidencia hexádico--octádica y funcional de Barbero--Immirzi}
\label{sec:barbero}
Sea $H\hookrightarrow O$ una inclusión marcada, $|H|=6$, $|O|=8$.
Los objetos de incidencia
\[
\mathcal F_6=H\times\binom H2,\qquad
\mathcal F_8=O\times\binom H2
\]
tienen cardinales $90$ y $120$. Su diferencia normalizada da
\[
\frac{|\mathcal F_6|-|\mathcal F_8|}{24\binom62}=-\frac1{12}.
\]
Esta es una realización combinatoria del mismo escalar.
Una igualdad numérica con \eqref{eq:minus12} no borra
la distinción entre los espacios sobre los que se ha calculado.

\begin{figure}[htbp]
\centering
\begin{tikzpicture}[>=Latex,every node/.style={font=\small}]
\draw[rounded corners,draw=ink] (-1.8,-1.35) rectangle (1.8,1.35);
\foreach \a in {0,60,...,300}{
\fill[ink] (\a:0.83) circle(2pt);}
\node at (0,-1.7) {$H:\ 6$ puntos};
\draw[->,thick,ink] (2.1,0)--(3.6,0);
\node[above] at (2.85,0.1) {inclusión marcada};
\draw[rounded corners,draw=ink] (4,-1.35) rectangle (7.6,1.35);
\foreach \a in {0,60,...,300}{
\fill[ink] ($(5.8,0)+(\a:0.83)$) circle(2pt);}
\fill[accent] (4.5,0) circle(2pt);
\fill[accent] (7.1,0) circle(2pt);
\node at (5.8,-1.7) {$O:\ 8$ puntos};
\node at (0,-2.35) {$6\binom62=90$};
\node at (5.8,-2.35) {$8\binom62=120$};
\end{tikzpicture}
\caption{Incidencia marcada de seis a ocho puntos. El dibujo representa
los conjuntos y la inclusión; no pretende convertir las caras o vértices
de un cubo en bloques de Steiner sin su mapa de incidencia.}
\end{figure}

\subsection{Los pesos preceden al coeficiente de polo}
En el módulo libre de sucesos orientados, la fuente fija la cadena
fundamental de la vuelta \cite{precedencia}
\[
c_{12}^{+}=
\sum_{j\in\Z/12\Z}[j,\mathcal F_6]
-[\partial_{\rm ret},\mathcal F_8].
\]
La multiplicidad doce corresponde a los sectores y la multiplicidad uno
a la costura orientada. La involución cambia la orientación de la cadena
completa. Sea
\[
S_m(q)=\frac{q^m}{1-q^{3m}},\qquad 0<q<1.
\]
El lector lineal asigna $S_{90}$ a cada sector y $S_{120}$ a la costura;
en consecuencia,
\[
\mathscr R(c_{12}^{+})=12S_{90}-S_{120}.
\]

\begin{theorem}[Coeficiente de polo y evaluación conjugada]
La imagen de la cadena fundamental tiene coeficiente $1/24$
respecto de $(1-q)^{-1}$ y residuo $-1/24$ respecto de $(q-1)^{-1}$.
Su evaluación en los canales previamente generados determina
\begin{equation}
\gamma_{\HMT}=
\frac{\pi AC^*}{180}
+12\bigl[S_{90}(q_-)-S_{90}(q_+)\bigr]
-\bigl[S_{120}(q_-)-S_{120}(q_+)\bigr].
\label{eq:gamma}
\end{equation}
\end{theorem}
\begin{proof}
La factorización $1-q^{3m}=(1-q)\sum_{j=0}^{3m-1}q^j$ implica
$\lim_{q\uparrow1}(1-q)S_m(q)=1/(3m)$. Por linealidad,
\[
\frac{12}{270}-\frac1{360}=\frac1{24}.
\]
La coordenada $q-1$ invierte el signo del residuo.
La evaluación conjugada y la componente angular de la fuente
producen \eqref{eq:gamma}.
\end{proof}

El certificado diofántico $4a-3b=45$ admite los pares
$(12+3t,1+4t)$, $t\ge0$; certifica la minimalidad de $(12,1)$,
pero no reemplaza el origen de las multiplicidades en la cadena.
La evaluación conservada en el propietario es
\[
\gamma_{\HMT}
=0.27400767269445927474725526077398041940\ldots.
\]
Esta cifra es una salida del funcional fuente; no es un valor externo
usado para seleccionar $A$, $C^*$ o sus coeficientes.

\subsection{Paridad y sensibilidad estructural}
Al invertir $C^*$ se intercambian $q_+$ y $q_-$ y cambia de signo
la componente angular. Por ello
\[
\gamma_{\HMT}(A,-C^*)=-\gamma_{\HMT}(A,C^*).
\]
El funcional conserva orientación, mientras que otros lectores,
como $q_+q_-$, son pares. Esta diferencia precede a su interpretación física.

\section{Elipse de acción y reciprocidad modular}
\label{sec:ellipse}
Fijadas las salidas $A,C^*,\hbar>0$ en la carta
$|C^*/A|<1$, definimos
\[
\eta_{\rm el}=\operatorname{artanh}(C^*/A),\quad
a_S=\sqrt{2\hbar}\,e^{\eta_{\rm el}/2},\quad
b_S=\sqrt{2\hbar}\,e^{-\eta_{\rm el}/2}.
\]
Las coordenadas canónicas equilibradas tienen unidad de raíz de acción;
no se suman aquí posición y momento sin una carta dimensional.
Se cumplen
\[
a_Sb_S=2\hbar,\qquad
\tau_{\rm el}=i\,\frac{b_S}{a_S},\qquad
R_{\rm el}=\sqrt{\frac{b_S}{a_S}}
=\left(\frac{A-C^*}{A+C^*}\right)^{1/4}.
\]
La raíz positiva determina la composición $\tau_{\rm el}=iR_{\rm el}^2$,
y permite recuperar la razón angular:
\[
\frac{C^*}{A}=\frac{1-R_{\rm el}^4}{1+R_{\rm el}^4}.
\]
La inversión de forma envía $R_{\rm el}$ a $R_{\rm el}^{-1}$ y
$\tau_{\rm el}$ a $-1/\tau_{\rm el}$ \cite{elipse}.

\begin{proposition}[Energía coincidente, acción distinguible]
Sea
\[
D(\eta)=\begin{pmatrix}e^\eta&0\\0&e^{-\eta}\end{pmatrix},
\quad
S_2=\begin{pmatrix}0&1\\1&0\end{pmatrix},
\quad
J_2=\begin{pmatrix}0&-1\\1&0\end{pmatrix}.
\]
Ambos mapas satisfacen $T^{\mathsf T}D(-\eta)T=D(\eta)$.
Sin embargo, para $\omega=\dd Q\wedge\dd P$,
\[
S_2^*\omega=-\omega,\qquad J_2^*\omega=\omega.
\]
Si $\mathcal I(\gamma)=\oint_\gamma P\,\dd Q$ sobre una curva cerrada,
\[
\mathcal I(S_2\gamma)=-\mathcal I(\gamma),\qquad
\mathcal I(J_2\gamma)=\mathcal I(\gamma).
\]
\end{proposition}
\begin{proof}
La conjugación intercambia las entradas diagonales de $D$.
Los determinantes de $S_2,J_2$ son $-1,1$, respectivamente, lo que
prueba la transformación de $\omega$. Para $\lambda=P\,\dd Q$,
\[
S_2^*\lambda=\dd(PQ)-\lambda,\qquad
J_2^*\lambda=\lambda-\dd(PQ).
\]
Integrar la diferencial exacta sobre la curva cerrada completa la prueba.
\end{proof}

La energía cuadrática y el módulo no bastan para distinguir ambos
transportes. La orientación de la acción sí los distingue, incluso
en la forma circular. El marcado geométrico es así parte de la lectura
recuperable, no un adorno añadido al valor.

\section{Operador de área y reconstrucción de la información sectorial}
\label{sec:area}
El reagrupamiento de seis trits construye una biyección entre
$(\Z/9\Z)^3$ y $(\Z/27\Z)^2$, ambos de cardinal $729$.
La biyección no se presenta como isomorfismo aditivo. Sobre un bloque
de borde $9\times9$, la fuente descompone 36 enlaces en dos
conjuntos de 18, etiquetados por los canales $90$ y $120$.

Sea $N_e=\operatorname{diag}(0,1,2)$ en cada enlace y
$N_{90},N_{120}$ las sumas por sector. Fijada la escala
\[
\theta_{\rm clk}=\frac{C^*}{6}\frac{\pi}{180},\qquad
a_0=\frac{1+2\cos(10^\circ)}3\theta_{\rm clk},
\]
el operador areal normalizado es
\[
\widehat{\mathcal A}/\ell_P^2
=\gamma_{\HMT}a_0(N_{90}+N_{120}).
\]
Su espectro tiene valores $n\gamma_{\HMT}a_0$, $0\le n\le72$,
con multiplicidades dadas por $(1+z+z^2)^{36}$.
La prueba consiste en sumar los autovalores $0,1,2$ de los
36 factores tensoriales \cite{area}.

\begin{theorem}[Recuperación de los sectores de borde]
Sean
\[
N=N_{90}+N_{120},\qquad D=90N_{90}+120N_{120}.
\]
Entonces
\begin{equation}
N_{90}=4N-\frac D{30},\qquad
N_{120}=\frac D{30}-3N.
\label{eq:recover}
\end{equation}
Por tanto, el total y el lector ponderado recuperan conjuntamente
las dos ocupaciones.
\end{theorem}
\begin{proof}
La matriz $\left(\begin{smallmatrix}1&1\\90&120\end{smallmatrix}\right)$
tiene determinante $30$. Su inversión da \eqref{eq:recover}.
En el dominio de ocupaciones enteras, la imagen satisface
$D\in30\Z$ y $90N\le D\le120N$.
\end{proof}

El Hamiltoniano modular de la fuente es
$K_A=-\log\rho_A=cI+\Delta_{\rm mod}D$.
Para $\Delta_{\rm mod}\ne0$ y normalización conocida,
$D=(K_A-cI)/\Delta_{\rm mod}$ publica el segundo observable.
Este operador describe la estructura espectral del estado reducido;
no se identifica por notación con el generador de evolución temporal.
El área conserva la ocupación total y la pareja añade la procedencia.
Por ejemplo, $(N_{90},N_{120})=(1,0)$ y $(0,1)$ comparten $N=1$,
pero tienen $D=90$ y $D=120$.

\input{sections/lectura_modular.tex}

\subsection{Estado reducido, descomposición de Schmidt y entropía de entrelazamiento}
Si $\rho_A=\sum_\nu p_\nu|\nu\rangle\langle\nu|$,
la purificación
\[
|\Psi\rangle=\sum_\nu\sqrt{p_\nu}\,
|\nu\rangle_A|\nu\rangle_B
\]
satisface $\Tr_B|\Psi\rangle\langle\Psi|=\rho_A$.
Así, $-\Tr(\rho_A\log\rho_A)$ es la entropía de entrelazamiento
de esta bipartición pura concreta. El estado ampliado declara
qué sistema porta esa interpretación.

Para una involución de dos hojas con un autovector de cada signo,
\[
\rho_y=\frac{e^{-yR}}{2\cosh y},\qquad R^2=I,
\]
la diagonalización proporciona
\[
S(\rho_y)=\log(2\cosh y)-y\tanh y,\qquad
D(\rho_y\Vert\rho_{-y})=2y\tanh y.
\]
La entropía es par en $y$; la distancia compara explícitamente
las dos orientaciones. No se identifica esta entropía de dos niveles
con la capacidad $81\log3$ del corte triádico.

\section{Acción, gravedad e información térmica}
\label{sec:gravity}
\subsection{Período dodecafásico y realización circular de la acción}
El calendario ya construido determina $T_{108}=108t_0$.
En su realización circular,
\[
\omega_{108}=\frac{2\pi}{108t_0},\qquad
R_{108}=\frac{c_{\rm int}}{\omega_{108}}
=\frac{54}{\pi}\ell_0,\qquad \ell_0=c_{\rm int}t_0.
\]
La unidad $\ell_0$ y las cartas de tiempo y velocidad se mantienen
explícitas; la carta SI no selecciona el ciclo.
Para cada sección orientada de acción $a\in\{\mathrm{pre},\mathrm{ret}\}$,
\[
E_{108,a}=\hbar_a\omega_{108},\qquad
m_{108,a}=\frac{E_{108,a}}{c_{\rm int}^2}.
\]
La longitud reducida y la completa son respectivamente
\[
\frac{\hbar_a}{m_{108,a}c_{\rm int}}=R_{108},
\qquad
\frac{h_a}{m_{108,a}c_{\rm int}}=2\pi R_{108}.
\]
La misma acción y la misma duración intervienen en ambas lecturas
\cite{autodual}.

\begin{proposition}[Selector de coincidencia de radios]
En la familia dimensional
\[
G_a(\vartheta)=\vartheta\frac{c_{\rm int}^5t_0^2}{\hbar_a},
\qquad \vartheta>0,
\]
la condición de la realización circular
$G_a(\vartheta)m_{108,a}/c_{\rm int}^2=R_{108}$
tiene la única solución
\[
\vartheta^\circ_{108}=(54/\pi)^2.
\]
\end{proposition}
\begin{proof}
El primer radio es
$r_g(\vartheta)=\vartheta(\pi/54)\ell_0$.
Compararlo con $(54/\pi)\ell_0$, con $\ell_0>0$, reduce la
igualdad a una ecuación lineal estrictamente creciente.
\end{proof}
La condición de coincidencia es parte explícita de esta realización;
el teorema resuelve su coeficiente sin usar un valor metrológico de $G$.
El radio aquí empleado es $Gm/c^2$, por lo que no se traslada
inadvertidamente el factor dos de la convención $2Gm/c^2$.

\begin{proposition}[Acción y gravedad recíprocas, área conservada]
En las dos cartas anteriores,
\[
\frac{G_{\rm pre}}{G_{\rm ret}}=R_{\rm act}^{-1},\qquad
\hbar_aG_a=\vartheta^\circ_{108}c_{\rm int}^5t_0^2,
\qquad
\ell_P^2=\frac{\hbar_aG_a}{c_{\rm int}^3}
=\vartheta^\circ_{108}\ell_0^2.
\]
\end{proposition}
\begin{proof}
La definición de $G_a$ contiene $\hbar_a^{-1}$ y el resto de factores
es común a las dos cartas. Tomar el cociente y el producto da
las tres igualdades.
\end{proof}
Esta reciprocidad explica por qué la información de orientación puede
variar sin que cambie el área. Se dispone así de una relación precisa
entre la elipse de acción, la escala de Planck y el operador areal
de la sección~\ref{sec:area}.

\subsection{Normalización térmica y energía por ciclo}
La relación térmica del corpus publica
\begin{equation}
k_B\Theta_{\rm clk}\log3
=\frac{h_a}{108t_0}
=\frac{\hbar_a}{t_P}=E_{P,a},
\qquad t_P=\frac{54}{\pi}t_0.
\label{eq:thermalclock}
\end{equation}
El factor $\log3$ mide la información de la terna equiprobable;
la energía del reloj incorpora la acción ya generada.
La identidad conserva el producto $k_B\Theta_{\rm clk}$ como sección
de energía y hace explícita la carta térmica utilizada \cite{termica}.

La realización radiativa proporciona, en esta escala,
\[
n_\gamma\ell_P^3=\frac{2\zeta(3)}{\pi^2\log^3 3},\qquad
\frac{u_\gamma\ell_P^3}{E_P}
=\frac{\pi^2}{15\log^4 3},\qquad
\frac{s_\gamma\ell_P^3}{k_B}
=\frac{4\pi^2}{45\log^3 3}.
\]
Aquí se encuentran dos invariantes distintos: $\log3$ conserva la
información TRIT y $\zeta(3)$ el momento de ocupación bosónica.
Su presencia en una misma expresión ofrece una conexión pertinente
con un artículo posterior sobre constantes y momentos espectrales,
sin convertir una constante en nombre intercambiable de la otra.

\begin{apunte}{Boltzmann y alcance editorial}
Conservar la ecuación \eqref{eq:thermalclock} en el cuerpo enlaza acción,
tiempo e información. Para exponer además el valor individual de $k_B$,
se reunirá aquí su carta de temperatura y el constructor de esa
normalización, en vez de deducir dos valores independientes de una
sola ecuación para su producto. Las fórmulas radiativas pueden quedar
como interludio o desarrollarse en otro artículo: el material íntegro
se preserva en la recopilación.
\end{apunte}

\subsection{Barbero--Immirzi en la realización Cartan--Holst}
Una vez construido $\gamma_{\HMT}$, la fuente lo realiza como
coeficiente de la acción de Holst. En la carta geométrica declarada,
sea $e^I$ una co-tétrada no degenerada, $\omega^{IJ}$ una conexión,
$\Sigma^{IJ}=e^I\wedge e^J$ y
$F^{IJ}=\dd\omega^{IJ}+\omega^I{}_K\wedge\omega^{KJ}$.
Con $\kappa=8\pi G/c^4$, la acción especializada es
\begin{equation}
\begin{aligned}
S[e,\omega,\Psi]
={}&\frac1{2\kappa c}\int
\Sigma_{IJ}\wedge(\star F^{IJ}+\gamma_{\HMT}^{-1}F^{IJ})\\
&-\frac{\Lambda}{\kappa c}\int\operatorname{vol}_e
+S_m[e,\omega,\Psi].
\end{aligned}
\label{eq:holst}
\end{equation}
La convención de corriente es
$\delta_\omega S_m=-\tfrac12\int\delta\omega^{IJ}\wedge\sigma_{IJ}$.
La variación de \eqref{eq:holst}, bajo las condiciones de borde
del propietario, da
\[
D_\omega[(\star+\gamma_{\HMT}^{-1}I)\Sigma^{IJ}]
=\kappa c\,\sigma^{IJ}.
\]
Se obtiene usando $\delta F=D_\omega\delta\omega$, integrando por partes
y haciendo nulo el coeficiente de la variación arbitraria
\cite{holst}. La orientación de la frontera y la normalización
de la corriente forman parte de esta fórmula.

En firma lorentziana, $\star^2=-I$ sobre bivectores. La identidad
\[
(\star+\gamma^{-1}I)^{-1}
=\frac{\gamma^2}{1+\gamma^2}(-\star+\gamma^{-1}I)
\]
se comprueba multiplicando ambos miembros: los términos cruzados se
cancelan y queda $(1+\gamma^{-2})I$ antes de normalizar.
Así, el parámetro ya construido interviene en la respuesta torsional,
no únicamente en un valor de tabla.
Las identidades
$D_\omega T^I=F^I{}_J\wedge e^J$ y $D_\omega F^{IJ}=0$
mantienen conectados torsión, curvatura y balance.

\begin{apunte}{gravedad como parte del argumento}
Esta sección permite discutir gravedad, área e información sin añadir
un capítulo desligado de la acción. El desarrollo definitivo reunirá
el transporte desde el registro TPK a la co-tétrada y la conexión,
sus clases de holonomía y las cartas de $G$ y $\Lambda$.
Se han conservado los propietarios completos de esa realización.
La métrica lorentziana, la carta de Minkowski, el radio recíproco y
las geometrías de curvatura constante deben colocarse donde sus
aplicaciones efectivas sean necesarias; el borrador no decide
todavía si sus desarrollos ocuparán el cuerpo o un suplemento.
\end{apunte}

\section{Transporte angular y continuación excepcional}
\label{sec:cp}
\label{sec:extension}
La relación con CKM debe exhibir un mapa, no sólo proximidad narrativa
con Barbero. En las coordenadas de la fuente,
\[
\symbf\theta=M(A,h_\theta,\Delta)^{\mathsf T},\quad h_\theta=C^*/6,
\qquad
M=\begin{pmatrix}
2&-5&28\\0&7&-25/2\\1&-20&-2/3
\end{pmatrix},
\quad \det M=-3857/6.
\]
Aquí $h_\theta$ es una coordenada angular, distinta de la constante de Planck $h$.
Sean $c_h=(-5,7,-20)^{\mathsf T}$ y
$\ell_h=(75,176,-150)/3857$. Se comprueba
$\ell_hc_h=1$ y que $\ell_h$ anula las otras columnas. Por tanto,
\[
\mathcal R_{\rm CKM}=I-2c_h\ell_h
=M\operatorname{diag}(1,-1,1)M^{-1}
\]
satisface $\mathcal R_{\rm CKM}^2=I$ y
$\mathcal R_{\rm CKM}M=M\operatorname{diag}(1,-1,1)$.
La prueba sigue de $(c_h\ell_h)^2=c_h\ell_h$.
Es el transporte de la inversión de hoja al espacio de mezcla.
La fórmula actúa sobre tres ángulos; en esta acción particular,
$\delta_{\rm CP}=9A$ permanece fija \cite{ckm}.
La operación completa exige conservar también el lector complejo de
amplitudes, que la fuente publica separadamente.
La procedencia de las columnas de $M$ se conserva con sus constructores
en la recopilación para el desarrollo autónomo posterior.

\subsection{Paridad de rutas y conjugación de amplitudes}
En la ruta $\rho=((i_t,j_t,\ell_t);d_t)$, los operadores del corpus
actúan como
\[
P:(i,j)\mapsto(i,-j),\qquad
C_{\rm hoja}:\Sigma\leftrightarrow\Pi,\quad C^*\mapsto-C^*,
\]
mientras $T$ invierte el orden del registro y cada dirección.
La acción discreta
\[
\mathcal A_\lambda(\rho)
=\sum_{t=2}^L\left(
1-\delta_{d_t,d_{t-1}}
+\lambda(1-\delta_{\ell_t,\ell_{t-1}})\right)
\]
se conserva bajo cada operación y bajo su composición.
La prueba es combinatoria: una permutación global del alfabeto preserva
la igualdad entre vecinos, y la inversión del orden conserva el número
de cambios. Esta identidad usa la ruta registrada y sus hojas
\cite{cpt}.

El lector complejo ya está explícito en el tratado. Con
$c_{ab}=\cos\theta_{ab}$ y $s_{ab}=\sin\theta_{ab}$,
\begin{equation}
\resizebox{\linewidth}{!}{$\displaystyle
V(\symbf\theta,\delta)=
\begin{pmatrix}
c_{12}c_{13}&s_{12}c_{13}&s_{13}e^{-i\delta}\\
-s_{12}c_{23}-c_{12}s_{23}s_{13}e^{i\delta}&
c_{12}c_{23}-s_{12}s_{23}s_{13}e^{i\delta}&s_{23}c_{13}\\
s_{12}s_{23}-c_{12}c_{23}s_{13}e^{i\delta}&
-c_{12}s_{23}-s_{12}c_{23}s_{13}e^{i\delta}&c_{23}c_{13}
\end{pmatrix}.$}
\label{eq:ckm-full}
\end{equation}
Los ángulos de la carta se convierten a radianes antes de aplicar
las funciones trigonométricas.

\begin{proposition}[Conjugación exacta del lector complejo]
Para ángulos reales,
\[
\overline{V(\symbf\theta,\delta)}
=V(\symbf\theta,-\delta),\qquad
J(\symbf\theta,-\delta)=-J(\symbf\theta,\delta),
\]
donde
$J=c_{12}c_{23}c_{13}^2s_{12}s_{23}s_{13}\sin\delta$.
\end{proposition}
\begin{proof}
En \eqref{eq:ckm-full}, todos los factores trigonométricos son reales.
La conjugación intercambia $e^{i\delta}$ y $e^{-i\delta}$; el seno
de la fase cambia de signo. Las dos identidades siguen término a término.
\end{proof}

Por tanto, la matriz $\mathcal R_{\rm CKM}$ no agota los operadores
que deben seguirse en la realización CP. El material reunido contiene
la paridad espacial de la ruta, la inversión de hoja y la conjugación
compleja de las amplitudes \cite{mezcla}. La conjugación antiunitaria
de la evolución tiene, a su vez, una residencia propia: si
$CJ_EC^{-1}=-J_E$ y $CHC^{-1}=H$, con los dominios preservados,
entonces $CU(t)C^{-1}=U(-t)$ para
$U(t)=\exp(-J_EHt/\hbar)$ \cite{conjugacion}.

\begin{apunte}{composición CP que debe quedar visible}
Se conservarán contiguos el transporte de rutas, el de marcos de
mezcla y el de amplitudes. Para un lector de realización $\mathcal V$,
la relación que expresa la compatibilidad es
$\mathcal V(CP_{\rm HMT}\rho)
=D_u\overline{\mathcal V(\rho)}D_d^\dagger$,
con las refasificaciones $D_u,D_d$ explícitas.
Las fuentes reunidas y las identidades precedentes permiten trabajar
ese cuadrado completo. Ni la fase fija de un mapa parcial ni una
declaración nominal de equivalencia sustituyen su comprobación.
Este apunte registra una tarea de composición; no cambia el teorema
de rutas ni el lector complejo ya publicados.
\end{apunte}

\subsection{Incidencia excepcional y continuación modular}

La rama de Witt comienza en un registro concreto, no en un grupo
excepcional postulado como generador. El soporte extraído de
$\mathbf K$ es una tétrada; el registro firmado determina una
hexada en el marco de incidencia declarado. Las estrellas y sus
prolongaciones conectan con los códigos, Leech y el álgebra
$V^\natural$. El grupo Monstruo actúa en esta última, no directamente
sobre las cifras de $\alpha$ \cite{witt}.

En el artículo definitivo se expondrán las aplicaciones que unan ese
corredor con el borde marcado usado en \eqref{eq:gamma}.
Moonshine, la invariancia modular y la dualidad de radio son
continuaciones relacionadas pero de tipos diferentes.
La igualdad de invariantes modulares no demuestra por sí sola una
equivalencia de todas las realizaciones físicas de cuerdas.
La extensión hasta teoría M se decidirá después de disponer de sus
mapas y de su residencia demostrativa íntegra.

\clearpage
\section{Conclusiones y desarrollo de la exposición}
El material reunido permite organizar el artículo mediante tres
lecturas complementarias. La actualización conserva memoria;
la acción distingue orientación y escala; el área y el Hamiltoniano
modular recuperan la procedencia de las ocupaciones.
El funcional de Barbero--Immirzi enlaza la incidencia
hexádico--octádica y la publicación angular dentro de esa organización.

La primera contribución de esta redacción consiste en reunir los
mapas que hacen precisa esa continuidad. El incremento $\eta_m$
se calcula desde el evento TPK, el escalar $-1/12$ conserva su
representación operatoria, y los pesos $12:-1$ preceden al
coeficiente de polo $1/24$. Las ecuaciones de acción se sitúan
antes de las magnitudes que las consumen.

La siguiente composición del manuscrito incorporará las pruebas
completas de los lectores regionales, las figuras del registro
y de Witt, la conexión de incidencia con la geometría areal
y las cartas dimensionales de reconocimiento. La recopilación
separada conserva ya sus propietarios sin reducirlos a las
fórmulas mostradas. El criterio de selección será la necesidad
matemática de cada antecedente, no una cifra de páginas.

\subsection{Definiciones de trabajo para la lectura estructural}
Las siguientes formulaciones reúnen conceptos ya presentes en las fuentes;
su función en el borrador es orientar la exposición, no atribuirles
una novedad histórica por cambiar su redacción.

\paragraph{Número genealógico.}
Objeto de lectura que retiene, además de la publicación escalar, la
clase de rutas compatibles, la hoja, la orientación y la memoria
pertinentes. El escalar constituye una proyección del objeto; la igualdad
de escalares se distingue de la equivalencia de estados transportados
\cite{segmento,registro}.

\paragraph{Sección orientada de acción.}
Elemento positivo de la recta de acción cuya elección de orientación
conserva la razón $R_{\rm act}$, y cuya realización sobre una curva
permite distinguir los transportes mediante la integral $\oint P\,\dd Q$.
La forma elíptica separa el área $2\pi\hbar$ de la deformación modular
\cite{accion,elipse}.

\paragraph{Funcional de incidencia orientada.}
La expresión $\Gamma_{6\to8}$ reúne la lectura angular bilineal y la
diferencia de los dos canales de incidencia. Sus pesos proceden de la
cadena dodecafásica y su signo conserva la inversión de hoja. Su
realización de Barbero--Immirzi incorpora este contenido a la respuesta
areal y torsional \cite{precedencia,holst}.

\paragraph{Recuperación de procedencia.}
Un conjunto de observables recupera una coordenada de origen cuando su
aplicación conjunta es invertible sobre el dominio declarado.
La pareja $(N,D)$ es un ejemplo explícito: recupera $N_{90},N_{120}$,
mientras $N$ por separado conserva sólo su suma. Esta propiedad local
no se extrapola sin prueba a todas las coordenadas del estado.

\begin{apunte}{jerarquía para la siguiente revisión}
Mantener como eje acción, orientación, incidencia y recuperación.
Las cifras comparativas acompañarán ese argumento. La ampliación
excepcional y las realizaciones geométricas se situarán por su
necesidad en la cadena, sin convertir la acumulación de nombres
en sustituto de las aplicaciones que los relacionan.
\end{apunte}

\clearpage
\appendix
\section{Notas de procedencia y alcance de la revisión}
La revisión conserva los resultados, las ecuaciones y las figuras del
prototipo anterior. Las adiciones reúnen la lectura geométrica, la
reversibilidad del sello y la normalización del Hamiltoniano modular.
Las formulaciones matemáticas se acompañan de sus hipótesis; las notas
de integración distinguen los desarrollos que se incorporarán a la
siguiente revisión.

La documentación de procedencia y los programas de comprobación se
mantienen en el suplemento digital. Sus identificadores de trabajo
no forman parte de la nomenclatura matemática del artículo. El
inventario técnico del prototipo se conserva íntegro en ese suplemento.

Las referencias que siguen corresponden a manuscritos de Rubén Ramos
Balsa y a desarrollos documentados del corpus HMT--MD, todavía inéditos.
Se citan para preservar atribución y procedencia, no como publicaciones
arbitradas. La versión destinada a lectura independiente incorporará
en el cuerpo las demostraciones previas que todavía se citan mediante
estos manuscritos.
\clearpage
\begin{thebibliography}{99}
\raggedright
\small
\setlength{\itemsep}{4pt}
\setlength{\parsep}{0pt}
\setlength{\parskip}{0pt}
\bibitem{geometria} R. Ramos Balsa. \emph{Álgebras cuadráticas y generadores concretos de los regímenes TRIT}.
Manuscrito inédito, sección de realizaciones elíptica, parabólica e hiperbólica.
Fuente conservada en el suplemento digital.
\bibitem{minkowski} R. Ramos Balsa. \emph{El radión: realización de Minkowski y holonomía orientada}.
Manuscrito inédito, sección de incidencia cúbica y orientación causal.
Fuente conservada en el suplemento digital.
\bibitem{segmento} Corpus HMT--MD. \emph{Construcción genealógica del continuo}.
Segmento marcado, carta positiva y continuación con acarreo. Manuscrito inédito.
\bibitem{app} Corpus HMT--MD. \emph{APP: soporte, operaciones y caminos}.
Fuente \texttt{03\_app\_ortograma.tex}; manuscrito inédito.
\bibitem{trit} Corpus HMT--MD. \emph{Estado local ternario y orientación del transporte}.
Fuente \texttt{02\_trit\_geometrias.tex}; manuscrito inédito.
\bibitem{tpk} Corpus HMT--MD. \emph{Árbol operatorio del TPK}.
Fuente \texttt{03h\_arbol\_operatorio\_tpk\_rev11.tex}; manuscrito inédito.
\bibitem{registro} Corpus HMT--MD. \emph{Registro dodecafásico firmado y reconstrucción integral del sello}.
Unidad U016; manuscrito inédito.
\bibitem{accion} Corpus HMT--MD. \emph{Derivación de la escala de acción}.
Fuentes \texttt{c28\_accion.tex} y delta c34, manuscritos inéditos.
\bibitem{retorno} Corpus HMT--MD. \emph{Carácter regional y mantisa reducida de acción; recta de acción}.
Deltas c35--c36, manuscritos inéditos.
\bibitem{precedencia} Corpus HMT--MD. \emph{Cadena fundamental dodecafásica y funcional hexádico--octádico}.
Fuentes c43, manuscritos inéditos.
\bibitem{area} Corpus HMT--MD. \emph{Operador de área, Hamiltoniano modular y memoria}.
Manuscritos inéditos.
\bibitem{ckm} Corpus HMT--MD. \emph{Transformación racional involutiva en el espacio de parámetros CKM}.
Fuente \texttt{c41\_ckm.tex}; manuscrito inédito.
\bibitem{witt} Corpus HMT--MD. \emph{Bandera y rama excepcional}.
Fuente \texttt{15\_bandera\_y\_rama\_excepcional.tex}; manuscrito inédito.
\bibitem{elipse} Coordinación HMT--MD, 6 de septiembre de 2026.
\emph{Forma elíptica, radio modular y acción orientada}. Manuscrito inédito.
\bibitem{eta} Coordinación HMT--MD, 6 de septiembre de 2026.
\emph{Actualización TPK, memoria publicada y determinantes dodecafásicos}.
Texto y verificador, manuscritos inéditos.
\bibitem{autodual} Corpus HMT--MD. \emph{Identidad autodual Planck--CT108}.
Realización circular y reciprocidad de las cartas de acción. Manuscritos inéditos.
\bibitem{termica} Corpus HMT--MD. \emph{Relación térmica intrínseca de la escala 108}.
Fuente c40; manuscrito inédito.
\bibitem{holst} Corpus HMT--MD. \emph{Realización Cartan--Holst--Barbero}.
Propietario B028, copia G03; sector gravitatorio completo en G02.
\bibitem{cpt} Corpus HMT--MD. \emph{Conjugación, paridad y reversión temporal en el registro}.
Fuente 08a; manuscrito inédito.
\bibitem{mezcla} Corpus HMT--MD. \emph{Mezcla CKM y violación CP}.
Fuente 74; manuscrito inédito.
\bibitem{conjugacion} Corpus HMT--MD. \emph{Conjugación antiunitaria y orientación temporal}.
Fuente 82; manuscrito inédito.
\end{thebibliography}
\end{document}
